\documentclass[a4paper,11pt]{article}

% ─── Packages ─────────────────────────────────────────────────────────────────
\usepackage[T1]{fontenc}
\usepackage[utf8]{inputenc}
\usepackage[english]{babel}
\usepackage{geometry}
\usepackage{graphicx}
\usepackage{hyperref}
\usepackage{xcolor}
\usepackage{booktabs}
\usepackage{longtable}
\usepackage{enumitem}
\usepackage{fancyhdr}
\usepackage{titlesec}
\usepackage{listings}
\usepackage{parskip}
\usepackage{microtype}
\usepackage{tcolorbox}
\usepackage{lmodern}
\usepackage{array}
\usepackage{tabularx}

% ─── Page geometry ────────────────────────────────────────────────────────────
\geometry{
  top    = 2.5cm,
  bottom = 2.5cm,
  left   = 2.5cm,
  right  = 2.5cm
}

% ─── Colours ──────────────────────────────────────────────────────────────────
\definecolor{canblue}{RGB}{30, 90, 160}
\definecolor{canlight}{RGB}{220, 235, 255}
\definecolor{cangray}{RGB}{245, 245, 245}
\definecolor{codegray}{RGB}{100, 100, 100}
\definecolor{notegreen}{RGB}{0, 120, 60}
\definecolor{warnorange}{RGB}{200, 100, 0}

% ─── Hyperref setup ───────────────────────────────────────────────────────────
\hypersetup{
  colorlinks   = true,
  linkcolor    = canblue,
  urlcolor     = canblue,
  citecolor    = canblue,
  pdftitle     = {CANgaroo User Manual},
  pdfauthor    = {CANgaroo Project},
  pdfsubject   = {CAN Bus Analyzer and Trace Tool},
  pdfkeywords  = {CAN, LIN, bus analyzer, automotive, trace}
}

% ─── Section styling ──────────────────────────────────────────────────────────
\titleformat{\section}
  {\Large\bfseries\color{canblue}}
  {\thesection}{1em}{}[\color{canblue}\titlerule]

\newcommand{\sectionbreak}{\clearpage}

\titleformat{\subsection}
  {\large\bfseries\color{canblue!80!black}}
  {\thesubsection}{1em}{}

\titleformat{\subsubsection}
  {\normalsize\bfseries}
  {\thesubsubsection}{1em}{}

% ─── Code listings ────────────────────────────────────────────────────────────
\lstset{
  basicstyle    = \ttfamily\small,
  keywordstyle  = \color{canblue}\bfseries,
  commentstyle  = \color{codegray}\itshape,
  stringstyle   = \color{notegreen},
  backgroundcolor = \color{cangray},
  frame         = single,
  rulecolor     = \color{canblue!30},
  breaklines    = true,
  showstringspaces = false,
  numbers       = left,
  numberstyle   = \tiny\color{codegray},
  numbersep     = 8pt,
  tabsize       = 4,
  captionpos    = b
}

% ─── tcolorbox styles ─────────────────────────────────────────────────────────
\tcbuselibrary{skins, breakable}

\newtcolorbox{notebox}[1][Note]{
  colback  = canlight,
  colframe = canblue,
  fonttitle = \bfseries,
  title    = #1,
  breakable
}

\newtcolorbox{warnbox}[1][Warning]{
  colback  = yellow!10,
  colframe = warnorange,
  fonttitle = \bfseries\color{warnorange},
  title    = #1,
  breakable
}

% ─── Header / Footer ──────────────────────────────────────────────────────────
\pagestyle{fancy}
\fancyhf{}
\fancyhead[L]{\small\color{canblue}\textbf{CANgaroo} User Manual}
\fancyhead[R]{\small\color{canblue}\nouppercase{\leftmark}}
\fancyfoot[C]{\small\thepage}
\renewcommand{\headrulewidth}{0.4pt}
\renewcommand{\headrule}{\color{canblue}\hrule width\headwidth height\headrulewidth}

% ─── Shorthand macros ─────────────────────────────────────────────────────────
\newcommand{\key}[1]{\texttt{\small[#1]}}
\newcommand{\menu}[1]{\textit{#1}}
\newcommand{\ui}[1]{\textbf{#1}}
\newcommand{\file}[1]{\texttt{#1}}

% ══════════════════════════════════════════════════════════════════════════════
\begin{document}

% ─── Title page ───────────────────────────────────────────────────────────────
\begin{titlepage}
  \centering
  \vspace*{1.5cm}

  \includegraphics[width=3.5cm]{../src/assets/cangaroo.png}\\[1.2em]

  {\Huge\bfseries\color{canblue} CANgaroo}\\[0.6em]
  {\Large\color{canblue!70!black} CAN Bus Analyzer \& Trace Tool}\\[3em]

  \textcolor{canblue}{\rule{0.6\textwidth}{1.5pt}}\\[3em]

  {\LARGE\bfseries User Manual}\\[1em]
  {\large Version 0.1}\\[4em]

  {\large Open-source automotive bus analysis software}\\[0.5em]
  {\large for CAN, CAN-FD, and LIN}\\[5em]

  \textcolor{canblue}{\rule{0.6\textwidth}{0.5pt}}\\[2em]
  {\normalsize \today}

  \vfill
  {\small This document describes CANgaroo and its features.\\
  CANgaroo is free and open-source software.\\[0.5em]
  \url{https://github.com/Schildkroet/CANgaroo}}
\end{titlepage}

% ─── Table of contents ────────────────────────────────────────────────────────
\tableofcontents
\newpage

% ══════════════════════════════════════════════════════════════════════════════
\section{Introduction}

CANgaroo is a free, open-source CAN bus analyzer and trace tool built with Qt6 and C++.
It provides a graphical interface for capturing, analyzing, decoding, generating, replaying,
and logging CAN, CAN-FD, and LIN bus traffic.

\subsection{Key Features}

\begin{itemize}[leftmargin=1.5em]
  \item Real-time message capture from multiple CAN/LIN interfaces simultaneously
  \item Aggregated and chronological trace views
  \item DBC and LDF database loading for signal decoding
  \item Live signal graph with time series, scatter, gauge, and text visualizations
  \item Manual and cyclic message transmission (TX Generator)
  \item File replay with speed control and channel mapping
  \item Python scripting engine for automation and custom logic
  \item Protocol decoders for UDS (ISO-TP) and J1939
  \item Export to candump, Vector ASC, MDF4, PCAP, and PCAP-NG
  \item Configurable workspaces saved as XML
  \item Light and Dark themes
\end{itemize}

\subsection{Supported Platforms}

\begin{table}[h]
  \centering
  \begin{tabularx}{0.8\textwidth}{lX}
    \toprule
    \textbf{Platform} & \textbf{Notes} \\
    \midrule
    Linux   & Full support including SocketCAN \\
    Windows & Full support; additional drivers for PEAK, Kvaser, Vector, TinyCAN \\
    \bottomrule
  \end{tabularx}
\end{table}

% ══════════════════════════════════════════════════════════════════════════════
\section{Installation and First Start}

\subsection{Pre-built Binaries}

Pre-built binaries for Linux and Windows are available on the GitHub releases page:\\
\url{https://github.com/Schildkroet/CANgaroo/releases}

Download the appropriate package for your platform and run the installer or extract
the archive. No compilation is required.

\subsection{Linux (Build from Source)}

To build from source, install the required Qt6 libraries and build tools:

\begin{lstlisting}[language=bash, caption={Build on Linux}]
sudo apt install qt6-base-dev qt6-serialport-dev qt6-serialbus-dev \
     qt6-charts-dev libnl-3-dev libnl-route-3-dev

git clone https://github.com/Schildkroet/CANgaroo.git
cd CANgaroo/src
qmake6
make -j$(nproc)
./cangaroo
\end{lstlisting}

\subsection{Windows (Build from Source)}

For a custom build with optional drivers:

\begin{lstlisting}[language=bash, caption={Optional driver flags}]
# Enable PEAK CAN driver
qmake6 CONFIG+=peakcan

# Enable Kvaser driver
qmake6 CONFIG+=kvaser

# Both optional drivers
qmake6 CONFIG+=peakcan CONFIG+=kvaser
\end{lstlisting}

\subsection{First Start}

On first launch, CANgaroo opens with an empty workspace and no interfaces configured.
The recommended first steps are:

\begin{enumerate}
  \item Open \menu{Measurement \textrightarrow{} Setup Interface} (\key{Ctrl+Alt+S}) to add
        your hardware interfaces.
  \item Optionally load a DBC or LDF database via the Setup dialog.
  \item Press \key{F5} or click \ui{Start} to begin capturing.
\end{enumerate}

% ══════════════════════════════════════════════════════════════════════════════
\section{User Interface Overview}

CANgaroo uses a tabbed, dockable multi-window layout. Each tab is an independent
workspace that can contain any combination of views arranged as docked panels.

\subsection{Main Window}

The main window consists of:

\begin{description}[leftmargin=2em, style=nextline]
  \item[Menu bar] Access to all application functions (File, Measurement, Trace, Window, Help).
  \item[Toolbar] Quick-access buttons: \ui{Start}, \ui{Stop}, \ui{Setup}, and \ui{Graph}.
  \item[Tab bar] Tabs, each containing an independent layout of views.
  \item[Status bar] Displays the application version.
\end{description}

\subsection{Creating and Arranging Views}

New views are created from the \menu{Window \textrightarrow{} New} submenu.
Available view types are listed in Section~\ref{sec:views}.
Views can be:

\begin{itemize}
  \item Docked inside a tab (drag the title bar to a docking zone).
  \item Floated as independent windows (drag outside the tab area).
  \item Closed individually (the \ui{X} button on the view title bar).
\end{itemize}

\subsection{Workspace Files}

The complete layout, all view configurations, database paths, filter states, and
TX generator messages are saved as an XML workspace file (\file{.xml}).

\begin{table}[h]
  \centering
  \begin{tabularx}{0.75\textwidth}{lX}
    \toprule
    \textbf{Action} & \textbf{Shortcut} \\
    \midrule
    New Workspace     & \key{Ctrl+N} \\
    Open Workspace    & \key{Ctrl+O} \\
    Save Workspace    & \key{Ctrl+S} \\
    Save As           & \key{Ctrl+Shift+S} \\
    \bottomrule
  \end{tabularx}
\end{table}

The \menu{File} menu also keeps a list of the eight most recently opened workspaces for
quick access.

% ══════════════════════════════════════════════════════════════════════════════
\section{Interface Setup}
\label{sec:setup}

\subsection{Setup Dialog}

Open the setup dialog via \menu{Measurement \textrightarrow{} Setup Interface} or
\key{Ctrl+Alt+S}.

The dialog shows a tree of \textbf{Networks} and the \textbf{Interfaces} within each network.
Right-click on any item to access context-menu actions.

\subsubsection{Adding a Network}

A network is a logical grouping of interfaces (e.g., powertrain network, body network).
Click \ui{Add Network} to create one. A network can contain multiple interfaces.

\subsubsection{Adding an Interface}

With a network selected, click \ui{Add Interface}. A dialog lists all interfaces
discovered by the installed drivers. Select the desired hardware channel and confirm.

\subsubsection{Database Files}

Each network can have one or more DBC (CAN) or LDF (LIN) database files associated
with it. These files provide the signal and message definitions needed for decoding.

\begin{itemize}
  \item Click \ui{Add Database} (or right-click \textrightarrow{} Add) to load a \file{.dbc}
        or \file{.ldf} file.
  \item \ui{Reload} re-parses the file from disk without restarting the measurement.
  \item \ui{Remove} unloads the database from the network.
\end{itemize}

\subsection{Supported Hardware Drivers}
\label{sec:drivers}

\begin{longtable}{p{3.2cm} p{2.2cm} p{7.8cm}}
  \toprule
  \textbf{Driver} & \textbf{Platform} & \textbf{Description} \\
  \midrule
  \endhead
  SocketCAN       & Linux    & Native Linux kernel CAN interface. Supports all network
                               devices exposed by the kernel (e.g.\ \texttt{can0},
                               \texttt{vcan0}). \\[0.4em]
  SLCAN           & All      & Serial Line CAN. Communicates with CAN adapters
                               that use the SLCAN ASCII protocol over a serial port. \\[0.4em]
  GrIP            & All      & Custom GrIP protocol adapter over serial or network. \\[0.4em]
  CandleAPI       & All      & USB CAN adapters based on the \textit{gs\_usb} protocol
                               (Geschwister Schneider USB/CAN adapters, CANtact, etc.). \\[0.4em]
  CANBlaster      & All      & Built-in virtual loopback / message generator for
                               development and testing. Enable via
                               \menu{Measurement \textrightarrow{} Driver \textrightarrow{} CANblaster}. \\[0.4em]
  Vector          & Windows  & Vector hardware (VN series etc.) via the Qt SerialBus
                               \texttt{vectorcan} plugin. \\[0.4em]
  Kvaser          & Windows  & Kvaser CAN adapters. Requires \texttt{CONFIG+=kvaser}
                               and the CANlib SDK. \\[0.4em]
  PEAK CAN        & Windows  & PCAN devices. Requires \texttt{CONFIG+=peakcan}
                               and the PCAN-Basic SDK. \\[0.4em]
  TinyCAN         & Windows  & TinyCAN adapters via Qt SerialBus. Enable via
                               \menu{Measurement \textrightarrow{} Driver \textrightarrow{} TinyCAN}. \\
  \bottomrule
\end{longtable}

\subsection{Starting and Stopping a Measurement}

\begin{itemize}
  \item \ui{Start} (\key{F5}): Opens all configured interfaces and begins capturing.
  \item \ui{Stop} (\key{Shift+F5}): Closes all interfaces and ends the capture session.
  \item \ui{Reload Interfaces} (\menu{Measurement \textrightarrow{} Reload Interfaces}):
        Re-scans all drivers for available hardware without restarting.
\end{itemize}

% ══════════════════════════════════════════════════════════════════════════════
\section{Views}
\label{sec:views}

\subsection{Trace View}

The Trace View is the primary display for captured bus traffic. It can be opened
multiple times (\key{Ctrl+Shift+T}) for independent filtering configurations.

\subsubsection{Display Modes}

\begin{description}[leftmargin=2em, style=nextline]
  \item[Aggregated mode] Shows one row per unique message ID. Each row is updated
        in-place with the latest data and a running occurrence counter. Ideal for
        monitoring message content.
  \item[Unified / Chronological mode] Shows every individual frame in the order it
        was received, appended at the bottom. Ideal for timing analysis.
\end{description}

Toggle between modes using the button in the Trace View toolbar.

\subsubsection{Categories}

The Trace View supports multiple category tabs:
\begin{itemize}
  \item \ui{Aggregated} — default CAN/LIN message view.
  \item \ui{UDS} — reassembled UDS (ISO-TP) messages decoded from raw CAN frames.
  \item \ui{J1939} — reassembled J1939 PGN messages.
\end{itemize}

\subsubsection{Filtering}

Each Trace View instance has independent filter settings:

\begin{description}[leftmargin=2em, style=nextline]
  \item[TX/RX toggle] Show or hide transmitted (TX) and/or received (RX) frames.
  \item[Message ID filter] Restrict display to a specific message ID or range.
  \item[Interface filter] Show only frames from a selected bus interface.
  \item[Advanced filter dialog] Combine multiple criteria using the filter dialog button.
\end{description}

\subsubsection{Timestamp Display}

Three timestamp modes are available via the column header or toolbar:
\begin{itemize}
  \item \textbf{Relative} — time since the start of the measurement.
  \item \textbf{Absolute} — wall-clock time.
  \item \textbf{Delta} — time since the previous frame of the same ID.
\end{itemize}

\subsubsection{Clearing the Trace}

Press \key{Esc} or use \menu{Trace \textrightarrow{} Clear} to clear all captured messages
from memory.

\subsection{Log View}

The Log View (\key{Ctrl+Shift+L}) displays application-level diagnostic messages
(driver errors, connection events, parsing warnings, etc.).

\begin{itemize}
  \item Click \ui{Clear} to empty the log.
  \item Use \ui{Export} to save the log contents to a file.
\end{itemize}

\subsection{Graph View}
\label{sec:graph}

The Graph View (\key{Ctrl+Shift+G}) plots decoded signal values in real time.
Multiple Graph Views can be open simultaneously.

\subsubsection{Adding Signals}

The left panel shows a signal tree derived from loaded DBC/LDF databases.
Use the search box to filter by signal name. Double-click a signal (or drag it into
the plot area) to add it to the current visualization.

\subsubsection{Visualization Types}

Each column in the Graph View can display one of the following types:

\begin{description}[leftmargin=2em, style=nextline]
  \item[Time Series] A scrolling line graph of signal value vs.\ time. Most
        commonly used for continuous signals.
  \item[Scatter Plot] XY scatter visualization for correlating two signal values.
  \item[Gauge] A circular dial showing the current value of a single signal,
        including configurable min/max range.
  \item[Text] Displays the latest decoded value as text, useful for enumerated
        or low-update-rate signals.
\end{description}

\subsubsection{Controls}

\begin{itemize}
  \item \ui{Zoom In / Out / Reset}: Scale the time axis.
  \item \ui{Time Window}: Set the scrolling window duration (e.g.\ show the last 10 s).
  \item Stale values are visually dimmed when no new frame has been received for a
        configurable period.
\end{itemize}

\subsubsection{Standalone Graph Window}

\key{Ctrl+Shift+B} opens a floating, undocked graph window that can be moved to a
second monitor.

\subsection{CAN Status View}

The CAN Status View shows live bus statistics for every active interface:

\begin{table}[h]
  \centering
  \begin{tabularx}{\textwidth}{lX}
    \toprule
    \textbf{Column} & \textbf{Description} \\
    \midrule
    Driver / Interface & Driver name and interface identifier \\
    State              & Bus state (Active, Warning, Passive, Bus-off) \\
    RX / TX Frames     & Total received / transmitted frame count \\
    RX / TX Errors     & Error frame counters \\
    RX Overrun         & Frames lost due to receive buffer overflow \\
    TX Dropped         & Frames dropped because the TX queue was full \\
    Bus Load           & Estimated bus utilisation in percent \\
    Warning / Passive / Bus-off & Escalating error state counters \\
    Restarts           & Number of automatic bus-off recoveries \\
    \bottomrule
  \end{tabularx}
\end{table}

Click \ui{Clear} to reset all counters. Statistics update automatically while a
measurement is running.

\subsection{Raw TX View (Message View)}
\label{sec:rawtx}

The Raw TX View allows composing and sending a single CAN or CAN-FD frame manually.

\subsubsection{Frame Fields}

\begin{description}[leftmargin=2em, style=nextline]
  \item[Message ID] Enter the frame ID in hexadecimal. Enable the \ui{Extended}
        checkbox for 29-bit IDs (EFF).
  \item[DLC] Data Length Code (0–8 for CAN, 0–64 for CAN-FD).
  \item[Data bytes] Edit individual bytes in the 8-column grid.
  \item[FD / BRS] Enable CAN-FD framing and Bit Rate Switch.
  \item[RTR] Set the Remote Transmission Request flag.
  \item[Interface] Select which hardware channel to transmit on.
\end{description}

\subsubsection{Signal-Based Editing}

When a DBC database is loaded and the entered ID matches a known message, a signal
table appears below the byte grid. Enter physical values into the signal rows;
CANgaroo encodes them into the frame bytes automatically. Multiplexed signals show
a mux selector dropdown.

Click \ui{Send} to transmit the composed frame once.

\subsection{TX Generator View}
\label{sec:txgen}

The TX Generator View manages a list of messages that are transmitted periodically
(cyclically) or on demand.

\subsubsection{Adding Messages}

\begin{itemize}
  \item \ui{Add from Database}: Browse loaded DBC messages and add one or more by
        name. The frame bytes and interval are pre-populated from the database.
  \item \ui{Add Manually}: Create a blank entry and fill in the ID, DLC, data, and
        interval manually.
\end{itemize}

\subsubsection{Message List Columns}

\begin{table}[h]
  \centering
  \begin{tabularx}{\textwidth}{lX}
    \toprule
    \textbf{Column} & \textbf{Description} \\
    \midrule
    Enable    & Checkbox to include/exclude this message from cyclic transmission \\
    ID        & CAN frame ID (hex) \\
    DLC       & Data length \\
    Data      & Frame payload (hex bytes), editable inline \\
    Interval  & Transmission interval in milliseconds \\
    Interface & Target hardware channel \\
    Count     & Number of times this message has been sent \\
    \bottomrule
  \end{tabularx}
\end{table}

\subsubsection{Toolbar Actions}

\begin{description}[leftmargin=2em, style=nextline]
  \item[Run All / Stop All] Start or stop cyclic transmission of all enabled messages.
  \item[Send Once] Transmit the selected message a single time immediately.
  \item[Enable / Disable All] Toggle the enable checkbox for every message in the list.
  \item[Random Payload] Fill the selected message's data bytes with random values
        (useful for stress testing).
  \item[Compact / Expanded Layout] Switch between a condensed list and an expanded
        view that shows the bit matrix editor.
  \item[Zoom Slider] Adjust the row height in the expanded layout.
\end{description}

\subsubsection{Bit Matrix Editor}

In expanded layout, each message row shows a bit matrix widget where individual bits
can be toggled by clicking. This is particularly useful for testing bit-level signal
encoding.

\subsection{Script View}
\label{sec:script}

The Script View provides an embedded Python 3 scripting engine that can interact
with the bus in real time. All bus access is available through the \texttt{cangaroo}
module which is automatically importable inside the engine.

\begin{notebox}
Python 3 must be installed and discoverable on the system \texttt{PATH} (Linux) or
bundled with the installation (Windows) for the Script View to function.
\end{notebox}

\subsubsection{Editor and Console}

The view is split into two panes:

\begin{description}[leftmargin=2em, style=nextline]
  \item[Editor (top)] Python code editor. Scripts can be typed directly, or loaded
        from and saved to a \file{.py} file using the \ui{Load} / \ui{Save} buttons.
  \item[Console (bottom)] Captures all \texttt{print()} output and exception
        tracebacks from the running script. Click \ui{Clear} to empty it.
\end{description}

The \ui{Auto-run} checkbox causes the script to start automatically each time a
measurement begins. \ui{Run} and \ui{Stop} control the script manually.

\subsubsection{The \texttt{cangaroo} Module}

Import the module at the top of every script:

\begin{lstlisting}[language=Python]
import cangaroo
\end{lstlisting}

The sections below describe every function and class in the module.

% ── Message class ─────────────────────────────────────────────────────────────
\subsubsection{Message Class}

\texttt{cangaroo.Message} represents a single CAN or LIN frame.

\begin{lstlisting}[language=Python, caption={Constructing and inspecting a Message}]
msg = cangaroo.Message()        # empty frame
msg = cangaroo.Message(0x123)   # frame with ID 0x123

# Properties (read/write unless noted)
msg.id        = 0x1A0          # CAN frame ID
msg.dlc       = 8              # data length (0-8 classic, 0-64 CAN-FD)
msg.extended  = True           # 29-bit extended ID (EFF)
msg.fd        = True           # CAN-FD frame
msg.brs       = True           # bit-rate switch (CAN-FD only)
msg.rtr       = False          # remote transmission request
msg.bustype   = "CAN"          # "CAN" or "LIN"

# Read-only properties
ts  = msg.timestamp            # float, seconds since measurement start
iid = msg.interface_id         # int, interface that produced the frame
rx  = msg.is_rx                # True if received, False if TX echo
\end{lstlisting}

\begin{lstlisting}[language=Python, caption={Accessing frame payload}]
msg.set_byte(0, 0xFF)          # write one byte
b = msg.get_byte(0)            # read one byte -> int

raw = msg.get_data()           # all bytes -> bytes object
msg.set_data(b'\xDE\xAD\x00') # set payload from bytes (updates dlc)
\end{lstlisting}

A LIN message is created with a dedicated helper:

\begin{lstlisting}[language=Python]
lin = cangaroo.make_lin_message(id=0x21, dlc=4)
\end{lstlisting}

% ── Send / Receive ────────────────────────────────────────────────────────────
\subsubsection{Sending and Receiving Frames}

\begin{lstlisting}[language=Python, caption={Send a single frame}]
msg = cangaroo.Message(0x100)
msg.dlc = 4
msg.set_data(b'\x01\x02\x03\x04')
cangaroo.send(msg)                     # sends on interface 0
cangaroo.send(msg, interface_id=1)     # sends on interface 1
\end{lstlisting}

\begin{lstlisting}[language=Python, caption={Blocking receive loop}]
while True:
    msgs = cangaroo.receive(timeout=1.0)   # blocks up to 1 second
    for msg in msgs:
        print(msg)
\end{lstlisting}

\texttt{receive()} returns a list (possibly empty on timeout) of \texttt{Message}
objects. By default transmitted frames are excluded; enable TX echo to include them:

\begin{lstlisting}[language=Python]
cangaroo.enable_tx_echo(True)
\end{lstlisting}

% ── RX filter ─────────────────────────────────────────────────────────────────
\subsubsection{RX Filtering}

An ID/mask filter limits which frames enter the \texttt{receive()} queue.
All other frames are silently discarded.

\begin{lstlisting}[language=Python, caption={Accept only ID 0x100 (standard frames)}]
cangaroo.set_filter(id=0x100, mask=0x7FF, extended=False)
\end{lstlisting}

\begin{lstlisting}[language=Python, caption={Accept a range of IDs (0x200-0x2FF)}]
cangaroo.set_filter(id=0x200, mask=0x700)
\end{lstlisting}

\begin{lstlisting}[language=Python, caption={Remove the filter (accept all frames)}]
cangaroo.clear_filter()
\end{lstlisting}

The \texttt{extended} argument constrains matching to standard (\texttt{False})
or extended (\texttt{True}) IDs. Omit it (or pass \texttt{None}) to match both.

% ── Periodic TX ───────────────────────────────────────────────────────────────
\subsubsection{Periodic (Cyclic) Transmission}

\begin{lstlisting}[language=Python, caption={Send a heartbeat every 100 ms}]
msg = cangaroo.Message(0x123)
msg.dlc = 2
msg.set_data(b'\xDE\xAD')

handle = cangaroo.send_periodic(msg, interval_ms=100)

# ... later ...
cangaroo.stop_periodic(handle)
\end{lstlisting}

Multiple periodic tasks can run simultaneously. Each returns an independent
integer handle. The task stops automatically when the script is stopped or the
measurement ends.

% ── Trace access ──────────────────────────────────────────────────────────────
\subsubsection{Trace Access}

\begin{lstlisting}[language=Python, caption={Read from the captured trace}]
n = cangaroo.trace_size()           # total number of stored messages

msgs = cangaroo.get_trace()         # get all messages
msgs = cangaroo.get_trace(count=50) # get last 50 messages

cangaroo.clear_trace()              # clear the trace buffer
\end{lstlisting}

% ── Measurement control ───────────────────────────────────────────────────────
\subsubsection{Measurement Control}

\begin{lstlisting}[language=Python]
if not cangaroo.measurement_running():
    cangaroo.start_measurement()

# ... do work ...

cangaroo.stop_measurement()
\end{lstlisting}

% ── Interface helpers ─────────────────────────────────────────────────────────
\subsubsection{Interface Helpers}

\begin{lstlisting}[language=Python, caption={List all active interfaces}]
for intf in cangaroo.interfaces():
    print(intf["id"], intf["name"])

name = cangaroo.interface_name(0)   # name of interface with ID 0
\end{lstlisting}

% ── Logging ───────────────────────────────────────────────────────────────────
\subsubsection{Logging}

Messages written via the logging functions appear in the CANgaroo Log View, not
just in the script console.

\begin{lstlisting}[language=Python]
cangaroo.log("script started")          # same as log_info
cangaroo.log_info("measurement OK")
cangaroo.log_warning("voltage low")
cangaroo.log_error("unexpected response")
\end{lstlisting}

% ── DBC / CAN database ────────────────────────────────────────────────────────
\subsubsection{CAN Database (DBC) Access}

When a DBC file is loaded in the Setup Dialog, the scripting engine can decode and
encode signals.

\paragraph{Listing databases}

\begin{lstlisting}[language=Python]
for db in cangaroo.databases():
    print(db["file"], db["network"])
    for msg in db["messages"]:
        print(" ", msg["name"], hex(msg["id"]), msg["signals"])
\end{lstlisting}

\paragraph{Decoding received frames}

\begin{lstlisting}[language=Python, caption={Decode signals from a received frame}]
msgs = cangaroo.receive()
for msg in msgs:
    result = cangaroo.decode(msg)
    if result:
        for name, sig in result["signals"].items():
            print(f"{name}: {sig['value']} {sig['unit']}")
\end{lstlisting}

The returned dictionary has the structure:

\begin{lstlisting}[language=Python]
{
    "message": "EngineData",   # DBC message name
    "id":      0x100,          # raw CAN ID
    "sender":  "ECU1",         # transmitting node (if defined)
    "signals": {
        "EngineSpeed": {
            "value":      3500.0,   # physical value
            "raw":        7000,     # raw integer value
            "unit":       "rpm",
            "min":        0.0,
            "max":        8000.0,
            "value_name": "",       # enum label if defined
        },
        ...
    }
}
\end{lstlisting}

\paragraph{Extracting a single signal}

\begin{lstlisting}[language=Python]
speed = cangaroo.signal_value(msg, "EngineSpeed")  # float or None
\end{lstlisting}

\paragraph{Finding a message definition}

\begin{lstlisting}[language=Python]
defn = cangaroo.find_message("EngineData")    # by name
defn = cangaroo.find_message(0x100)           # by raw ID
defn = cangaroo.lookup(received_msg)          # by live BusMessage
# defn is None if the ID is not in any loaded DBC
\end{lstlisting}

\paragraph{Encoding frames from signal values}

\begin{lstlisting}[language=Python, caption={Build and send a frame by signal name}]
msg = cangaroo.encode("EngineData", {
    "EngineSpeed": 3500.0,
    "EngineTemp":  90.0,
})
cangaroo.send(msg)
\end{lstlisting}

\texttt{encode()} accepts either the message name (string) or raw CAN ID (integer)
as the first argument.

% ── LIN database ─────────────────────────────────────────────────────────────
\subsubsection{LIN Database (LDF) Access}

\begin{lstlisting}[language=Python, caption={Decode a received LIN frame}]
msgs = cangaroo.receive()
for msg in msgs:
    if msg.bustype == "LIN":
        result = cangaroo.decode_lin(msg)
        if result:
            for name, sig in result["signals"].items():
                print(f"{name}: {sig['value']} {sig['unit']}")
\end{lstlisting}

\begin{lstlisting}[language=Python, caption={Look up a LIN frame definition}]
frame = cangaroo.find_lin_frame("WheelSpeed")   # by name
frame = cangaroo.find_lin_frame(0x21)           # by frame ID
# frame["signals"] is a list of signal dicts
\end{lstlisting}

\begin{lstlisting}[language=Python, caption={List all loaded LDF databases}]
for db in cangaroo.lin_databases():
    print(db["file"], db["master"], db["speed"])
    for f in db["frames"]:
        print(" ", f["name"], f["id"], f["length"])
\end{lstlisting}

% ── AUTOSAR E2E ───────────────────────────────────────────────────────────────
\subsubsection{AUTOSAR E2E Profile 2}

CANgaroo includes built-in support for computing AUTOSAR End-to-End Protection
Profile~2 checksums (CRC-8H2F, polynomial 0x2F).

The protected header occupies the first two bytes of the frame:

\begin{table}[h]
  \centering
  \begin{tabularx}{0.7\textwidth}{llX}
    \toprule
    \textbf{Byte} & \textbf{Bits} & \textbf{Content} \\
    \midrule
    0 & 7:0 & CRC-8H2F checksum \\
    1 & 3:0 & Counter (0–14) \\
    1 & 7:4 & Reserved (zero) \\
    \bottomrule
  \end{tabularx}
\end{table}

\begin{lstlisting}[language=Python, caption={One-shot protect and send}]
DATA_ID = 0x1234
counter = 0

msg = cangaroo.encode("EngineData", {"EngineSpeed": 3500.0})
cangaroo.e2e_p2_protect(msg, data_id=DATA_ID, counter=counter)
cangaroo.send(msg)
\end{lstlisting}

\texttt{e2e\_p2\_protect()} writes the counter nibble into byte~1 and the computed
CRC into byte~0 in-place. The counter should be incremented (modulo~15) with each
transmission.

\begin{lstlisting}[language=Python, caption={Manual CRC placement}]
msg.set_byte(1, counter & 0x0F)                          # write counter first
crc = cangaroo.e2e_p2_compute_crc(msg, data_id=DATA_ID)  # compute CRC
msg.set_byte(0, crc)                                     # place CRC
cangaroo.send(msg)
\end{lstlisting}

\texttt{e2e\_p2\_compute\_crc()} reads byte~0 as \texttt{0x00} regardless of its
current content and returns the CRC byte without modifying the message.

% ── Complete example ──────────────────────────────────────────────────────────
\subsubsection{Complete Script Example}

The following script monitors the bus, decodes every received frame that has a
DBC definition, and forwards an E2E-protected response.

\begin{lstlisting}[language=Python, caption={Monitor, decode, and respond with E2E protection}]
import cangaroo

DATA_ID = 0x0042
counter = 0

cangaroo.set_filter(id=0x100, mask=0x7FF)   # listen only to ID 0x100

while True:
    msgs = cangaroo.receive(timeout=2.0)
    if not msgs:
        cangaroo.log_warning("timeout - no frames received")
        continue

    for msg in msgs:
        result = cangaroo.decode(msg)
        if result is None:
            continue

        speed = result["signals"].get("EngineSpeed", {}).get("value")
        if speed is None:
            continue

        cangaroo.log(f"EngineSpeed = {speed} rpm")

        # Build and send a response with E2E protection
        resp = cangaroo.encode("StatusReply", {"Speed": speed, "Status": 1.0})
        cangaroo.e2e_p2_protect(resp, data_id=DATA_ID, counter=counter)
        cangaroo.send(resp)
        counter = (counter + 1) % 15       # counter wraps at 15
\end{lstlisting}

\subsection{Replay View}
\label{sec:replay}

The Replay View loads a previously saved trace file and replays it through the
active hardware interfaces (or the virtual loopback).

\subsubsection{Supported Input Formats}

\begin{itemize}
  \item candump log (\file{.log})
  \item Vector ASC (\file{.asc})
  \item PCAP (\file{.pcap})
  \item PCAP-NG (\file{.pcapng})
\end{itemize}

\subsubsection{Playback Controls}

\begin{description}[leftmargin=2em, style=nextline]
  \item[Play / Pause / Stop] Standard transport controls.
  \item[Seek slider] Drag to jump to any position in the file. The current
        timestamp is shown next to the slider.
  \item[Speed multiplier] Increase playback speed (e.g.\ 2× replays at twice
        the original rate). Values less than 1× slow down playback.
  \item[Loop] When checked, the file restarts automatically after the last frame.
  \item[Auto-play] When checked, playback begins immediately after a file is loaded.
\end{description}

\subsubsection{Message Filter}

A tree on the left side of the view lists all interfaces and directions (TX/RX)
found in the file. Uncheck any entry to suppress those frames during replay.
\ui{Select All} and \ui{Deselect All} buttons allow bulk toggling.

\subsubsection{Channel Mapping}

When the replay file was recorded on different interface names than those currently
active, use the channel mapping table to redirect each recorded channel to an
available hardware channel.

% ══════════════════════════════════════════════════════════════════════════════
\section{Database Support}

\subsection{DBC Files (CAN)}

DBC (\textit{Database CAN}) is the standard format used by Vector tools to describe
CAN networks. CANgaroo's built-in parser supports:

\begin{itemize}
  \item Message definitions (ID, DLC, transmitting node)
  \item Signal definitions (bit position, length, byte order, value type)
  \item Signal scaling (factor, offset, min, max, unit)
  \item Value descriptions (enumeration labels)
  \item Multiplexed signals with multiplexer indicator
  \item Node definitions
  \item Message and signal attributes
\end{itemize}

\subsection{LDF Files (LIN)}

LDF (\textit{LIN Description File}) describes LIN bus networks. The parser supports:

\begin{itemize}
  \item Frame definitions and their signals
  \item Signal bit positions and lengths
  \item Schedule tables
  \item Node definitions (master, slaves)
  \item LIN enhanced checksum calculation for correct ASC export
\end{itemize}

\subsection{Loading Databases}

Databases are loaded per network in the Setup Dialog (Section~\ref{sec:setup}).
Once loaded, their message and signal definitions are available in all views:

\begin{itemize}
  \item Trace View shows the decoded message name next to the ID.
  \item Graph View populates the signal tree.
  \item Raw TX View and TX Generator View enable signal-based frame editing.
\end{itemize}

% ══════════════════════════════════════════════════════════════════════════════
\section{Protocol Decoders}

\subsection{UDS / ISO-TP Decoder}

The UDS decoder reassembles multi-frame ISO-TP (ISO 15765-2) messages and
interprets the payload as UDS (ISO 14229) services.

\begin{itemize}
  \item Supports single frames, first frames, consecutive frames, and flow control.
  \item Interprets service identifiers (e.g.\ 0x22 ReadDataByIdentifier,
        0x2E WriteDataByIdentifier, 0x27 SecurityAccess).
  \item Decodes Negative Response Codes (NRC) for service 0x7F.
  \item Reassembled messages appear in the \ui{UDS} category tab of the Trace View.
\end{itemize}

\subsection{J1939 Decoder}

The J1939 decoder handles the higher-layer protocol used in commercial vehicles.

\begin{itemize}
  \item Extracts PGN, Source Address (SA), and Destination Address (DA) from
        29-bit extended CAN IDs.
  \item Reassembles multi-packet messages using the Transport Protocol:
    \begin{itemize}
      \item BAM (Broadcast Announce Message) sessions
      \item CM (Connection Management) point-to-point sessions
    \end{itemize}
  \item Reassembled messages appear in the \ui{J1939} category tab of the Trace View.
\end{itemize}

% ══════════════════════════════════════════════════════════════════════════════
\section{Saving and Exporting Traces}

\subsection{Saving During Measurement}

Use \menu{Trace \textrightarrow{} Save Trace to File} to write all currently buffered
messages to a file. A format selection dialog lets you choose the output format.

\subsection{Supported Export Formats}

\begin{longtable}{p{3cm} p{2.5cm} p{7.5cm}}
  \toprule
  \textbf{Format} & \textbf{Extension} & \textbf{Notes} \\
  \midrule
  \endhead
  candump log     & \file{.log}     & Plain-text format compatible with the Linux
                                      \texttt{can-utils} \texttt{candump} tool.
                                      Supports CAN-FD (\texttt{\#\#flag}) and
                                      RTR (\texttt{\#Rn}) notation. \\[0.4em]
  Vector ASC      & \file{.asc}     & Standard automotive logging format.
                                      Compatible with Vector CANalyzer/CANoe and
                                      other tools. Includes date/time header and
                                      LIN enhanced checksum. \\[0.4em]
  MDF4            & \file{.mdf}     & Binary MDF4 format with structured record
                                      layout. Supports up to 64-byte CAN-FD payloads.
                                      Compatible with analysis tools that read MDF4. \\[0.4em]
  PCAP            & \file{.pcap}    & Wireshark-compatible packet capture format. \\[0.4em]
  PCAP-NG         & \file{.pcapng}  & Extended PCAP-NG with per-interface metadata
                                      blocks. Modern Wireshark format. \\
  \bottomrule
\end{longtable}

\subsection{Full Trace Export and Import}

\menu{Trace \textrightarrow{} Export Full Trace} and \menu{Trace \textrightarrow{} Import Full Trace}
allow saving and restoring the entire in-memory trace buffer, including messages
that have scrolled off the view.

% ══════════════════════════════════════════════════════════════════════════════
\section{Conditional Logging}

The Conditional Logging feature writes frames to a file only when a user-defined
signal condition is met.

\begin{enumerate}
  \item Open \menu{Measurement \textrightarrow{} Conditional Logging}.
  \item Specify an output file path.
  \item Add one or more conditions using the signal tree browser:
    \begin{itemize}
      \item Select a signal from a loaded DBC database.
      \item Choose a comparison operator (\texttt{==}, \texttt{!=}, \texttt{<},
            \texttt{>}, \texttt{<=}, \texttt{>=}).
      \item Enter the threshold value.
    \end{itemize}
  \item Select whether multiple conditions must \textit{all} be true (\ui{AND})
        or \textit{any one} must be true (\ui{OR}).
  \item Enable logging with the \ui{Enable} checkbox.
\end{enumerate}

When the condition evaluates to true, the current frame is appended to the output
file in the selected format.

% ══════════════════════════════════════════════════════════════════════════════
\section{Application Settings}

Open \menu{File \textrightarrow{} Settings} (\key{Ctrl+Alt+S}) to adjust
application-wide preferences.

\begin{table}[h]
  \centering
  \begin{tabularx}{\textwidth}{lX}
    \toprule
    \textbf{Setting} & \textbf{Description} \\
    \midrule
    Theme            & Switch between \ui{Light} and \ui{Dark} visual themes \\
    Language         & Select the UI language (restart may be required) \\
    Restore Windows  & Restore dock layout and window geometry on next launch \\
    Clear on Start   & Automatically clear the trace when a measurement starts \\
    Font Size        & Adjust the application-wide font point size \\
    \bottomrule
  \end{tabularx}
\end{table}

% ══════════════════════════════════════════════════════════════════════════════
\section{Keyboard Shortcuts Reference}

\begin{longtable}{p{4.5cm} p{8.5cm}}
  \toprule
  \textbf{Shortcut} & \textbf{Action} \\
  \midrule
  \endhead
  \key{F5}               & Start Measurement \\
  \key{Shift+F5}         & Stop Measurement \\
  \key{Ctrl+N}           & New Workspace \\
  \key{Ctrl+O}           & Open Workspace \\
  \key{Ctrl+S}           & Save Workspace \\
  \key{Ctrl+Shift+S}     & Save Workspace As \\
  \key{Ctrl+Alt+S}       & Setup Interface / Settings \\
  \key{Ctrl+Shift+T}     & New Trace View \\
  \key{Ctrl+Shift+L}     & New Log View \\
  \key{Ctrl+Shift+G}     & New Graph View \\
  \key{Ctrl+Shift+B}     & New Standalone Graph Window \\
  \key{Esc}              & Clear Trace \\
  \key{Alt+F4}           & Exit Application \\
  \bottomrule
\end{longtable}

% ══════════════════════════════════════════════════════════════════════════════
\section{Typical Workflows}

\subsection{Monitor a CAN Bus}

\begin{enumerate}
  \item Connect your CAN adapter.
  \item Configure the interface in \menu{Measurement \textrightarrow{} Setup Interface}.
  \item Optionally load a DBC file for the target network.
  \item Press \key{F5} to start.
  \item The Trace View immediately shows incoming frames. Switch to
        \ui{Aggregated} mode to see only the latest value of each ID,
        or \ui{Unified} mode for the full chronological log.
\end{enumerate}

\subsection{Analyze Signals with the Graph View}

\begin{enumerate}
  \item Complete the steps above with a DBC database loaded.
  \item Open a Graph View (\key{Ctrl+Shift+G}).
  \item In the signal tree, search for the signal of interest.
  \item Double-click it to add a Time Series plot.
  \item Adjust the time window and zoom to inspect the waveform.
\end{enumerate}

\subsection{Transmit Cyclic Messages}

\begin{enumerate}
  \item Open a TX Generator View (\menu{Window \textrightarrow{} New \textrightarrow{} Generator View}).
  \item Click \ui{Add from Database} and select the message to transmit.
  \item Set the desired interval (ms) and edit the payload if needed.
  \item Enable the message and click \ui{Run All}.
\end{enumerate}

\subsection{Replay a Saved Log File}

\begin{enumerate}
  \item Open a Replay View (\menu{Window \textrightarrow{} New \textrightarrow{} Replay View}).
  \item Click the folder icon or \ui{Open File} and select your log file.
  \item Optionally adjust the speed multiplier and filter the channels.
  \item Click \ui{Play}.
\end{enumerate}

\subsection{Automate with Python}

\begin{enumerate}
  \item Open a Script View (\menu{Window \textrightarrow{} New \textrightarrow{} Script View}).
  \item Write or load your Python script.
  \item Press \ui{Run}. The console pane shows output and any errors.
  \item Enable \ui{Auto-run} to have the script start automatically with each
        measurement.
\end{enumerate}

% ══════════════════════════════════════════════════════════════════════════════
\section{Troubleshooting}

\begin{description}[leftmargin=2em, style=nextline]

  \item[Interface not appearing in Setup Dialog]
    Click \menu{Measurement \textrightarrow{} Reload Interfaces}. On Linux,
    ensure the CAN interface is up (\texttt{ip link set can0 up type can bitrate 500000}).
    Check the Log View for driver error messages.

  \item[No messages received]
    Verify the bitrate matches the network. On Linux, confirm the interface state
    with \texttt{ip -details link show can0}. Check the CAN Status View for
    bus-off or error-passive conditions.

  \item[DBC signals not decoded]
    Confirm the DBC file is loaded in the Setup Dialog for the correct network.
    Ensure the message ID in the DBC matches the ID on the bus (standard vs.\
    extended frame type).

  \item[Graph View shows no signals]
    A DBC database must be loaded and the measurement must be running. The signal
    tree is populated from the database; values appear only when matching frames
    arrive.

  \item[Python script does not start]
    Verify Python 3 is installed and on the system \texttt{PATH}. Check the
    console pane for import errors or syntax exceptions.

  \item[Replay file not loading]
    Ensure the file extension matches one of the supported formats (\file{.log},
    \file{.asc}, \file{.pcap}, \file{.pcapng}). Check the Log View for parsing
    error messages.

  \item[Application crashes on start]
    Delete the workspace XML file and restart. If the crash persists, launch from
    a terminal to capture the error output.

\end{description}

% ══════════════════════════════════════════════════════════════════════════════
\section{License and Contributing}

CANgaroo is released under the \textbf{GNU General Public License v2.0}.
You are free to use, modify, and distribute it under the terms of that license.

The source code is hosted on GitHub:
\url{https://github.com/Schildkroet/CANgaroo}

Bug reports and feature requests are welcome via the GitHub issue tracker.
Pull requests should follow these coding conventions:

\begin{itemize}
  \item Modern C++20, Qt6 API only
  \item Allman brace style
  \item \texttt{override} instead of \texttt{Q\_DECL\_OVERRIDE}
  \item Scoped enums (\texttt{enum class})
  \item Smart pointers over raw owning pointers
\end{itemize}

% ──────────────────────────────────────────────────────────────────────────────
\vfill
\begin{center}
  \textcolor{canblue}{\rule{0.5\textwidth}{0.5pt}}\\[0.5em]
  {\small CANgaroo User Manual \quad|\quad \today}
\end{center}

\end{document}
